\section{Keeping the partition random}
\label{l:random-partition-section}

Selecting one partition and imposing common coset caps already handles the
averaged marginal hypothesis. Here we construct a different retained law:
each retained permutation is good for many partitions, without choosing one
partition on which the law must satisfy every cap. The partition therefore
varies inside the coefficient estimate. Its event of simultaneous goodness
must be inserted before decomposing a test vector.

\begin{proposition}[Random-partition lift]\label{l:random-partition}
Let $n$ tend to infinity through multiples of a fixed integer $q>3$,
and let $\mu$ be probability laws on $S_n$.  For a uniform
$n/q$-subset $C$, put $H_C=\Sym(C)$, fixing its complement pointwise.
Suppose
\[
 \mathbb E_C\|\mu_{H_C}-(U_G)_{H_C}\|_{\TV}=o(1),
 \qquad \widehat\mu(\sgn)=0.
\]
There is a conditional probability law $\nu$ with
$\|\nu-\mu\|_{\TV}=o(1)$ such that, for every irreducible $\rho$
of degree $D$,
\begin{equation}\label{l:random-op}
 \|\widehat\nu(\rho)\|_{\op}
 \le \frac2M\sqrt{2qC_1}\,D^{-(1-3/q)/2},
\end{equation}
where $M=1-o(1)$ is the retained mass and $C_1$ is from
Lemma~\ref{l:degree-basic}.  In particular for $q=8$,
$\|\mu^{*8}-U_G\|_{\TV}=o(1)$.
\end{proposition}
\begin{proof}
Put $h_0=|H_C|/|G|$.  Say $g$ is good for $C$ if
$\mu(gH_C)\le2h_0$.  The total $\mu$-mass of bad cosets is at most
twice their marginal total-variation distance.  Thus the set
\[
 R=\{g:\mathbb P_C(g\text{ is bad for }C)\le1/(2q)\}
\]
has mass $M=1-o(1)$ by Markov's inequality.  For every $g\in R$ a
uniform ordered equal partition $(C_1',\ldots,C_q')$ has probability
at least $1/2$ that $g$ is good for all blocks.  Define
$\nu=\mu(\cdot\mid R)$.

We use Lemma~\ref{l:coefficient-evaluation}: if a function $F$ on
a finite group is in the span of matrix coefficients of a collection
$\Lambda$ of irreducibles, then
\begin{equation}\label{l:random-evaluation}
 \max|F|^2\le
 \left(\sum_{\tau\in\Lambda}(\dim\tau)^2\right)
                     \mathbb E_U|F|^2.
\end{equation}
The sum counts distinct types, since equivalent copies have the same
space of coefficient functions.

Fix a partition and restrict $\rho$ to its direct product of block
subgroups.  Decompose orthogonally into tensor-product constituents,
including all multiplicity spaces.  Every constituent has product
of factor degrees at most $D$.  Assign it to one factor with degree
at most $D^{1/q}$.  This splits a vector $v=\sum_i v_i$
orthogonally, with $v_i$ using only $H_{C_i'}$-types of degrees at
most $D^{1/q}$.  Fix a test vector $w$ and set
$F_i(g)=\langle w,\rho(g)v_i\rangle$.  On a coset $gH_{C_i'}$
it belongs to the coefficient space in
\eqref{l:random-evaluation}, whose dimension is at most
\[
 \sum_{b\le D^{1/q}}b^2\le C_1D^{3/q}.
\]
The sum counts each irreducible type once; multiplicities do not
increase the function space of matrix coefficients.  Summing the
coset maximum bound only over good cosets, and using Schur
orthogonality on $G$, yields
\begin{equation}\label{l:good-coset-square}
 \sum_{g\text{ good for }C_i'}\mu(g)|F_i(g)|^2
 \le2C_1D^{3/q-1}\|w\|^2\|v_i\|^2.
\end{equation}
Only a bound on total coset mass was used here.

For $g\in R$ insert the all-good partition event, whose probability
is at least $1/2$, before splitting the coefficient.  It follows that
\[
 |\langle w,\widehat\nu(\rho)v\rangle|
 \le\frac2M\mathbb E_{(C_i')}
      \sum_i\sum_{g\text{ good for }C_i'}\mu(g)|F_i(g)|.
\]
Cauchy--Schwarz and \eqref{l:good-coset-square}, followed by
$\sum_i\|v_i\|\le\sqrt q\|v\|$, prove
\eqref{l:random-op}.  The dependence of $v_i$ on the partition
is harmless because the estimate holds for each fixed partition.

The sign coefficient obeys
$|\widehat\nu(\sgn)|\le(1-M)/M=o(1)$, since the original sign
coefficient is zero. The same argument works when that coefficient is $o(1)$.  Fourier Plancherel in the normalization of
Lemma~\ref{l:plancherel} gives
\[
 4\|\nu^{*r}-U_G\|_{\TV}^2
 \le\sum_{\rho\ne\mathbf1}D_\rho
                    \|\widehat\nu(\rho)^r\|_{\HS}^2.
\]
For $q=r=8$, the nonsign sum is at most an absolute constant times
\[
 \sum_{\rho\ne\mathbf1,\sgn}
           D_\rho^{2-r(1-3/q)}
 =\sum_{\rho\ne\mathbf1,\sgn}D_\rho^{-3}=o(1)
\]
by Lemma~\ref{l:degree-basic}.  The sign term also vanishes, and
$\|\mu^{*8}-\nu^{*8}\|_{\TV}\le8(1-M)=o(1)$.
\end{proof}
